\documentclass{article}

\usepackage{amsmath,amssymb}
\usepackage{xurl}
\usepackage[hidelinks]{hyperref}
\setlength{\emergencystretch}{2em}

% Shared commands for notes in docs/paper-gaps/.

\DeclareUnicodeCharacter{2080}{\ifmmode{}_{0}\else\textsubscript{0}\fi}
\DeclareUnicodeCharacter{2081}{\ifmmode{}_{1}\else\textsubscript{1}\fi}
\DeclareUnicodeCharacter{2082}{\ifmmode{}_{2}\else\textsubscript{2}\fi}
\DeclareUnicodeCharacter{2099}{\ifmmode{}_{n}\else\textsubscript{n}\fi}

\newcommand{\Lean}{\textsc{Lean}}
\ExplSyntaxOn
\NewDocumentCommand{\leanid}{m}{
  \ifmmode
    \textsf{\detokenize{#1}}
  \else
    \group_begin:
    \urlstyle{sf}
    \tl_set:Nn \l_tmpa_tl {#1}
    \tl_replace_all:Nnn \l_tmpa_tl {\_} {_}
    \exp_args:NV \path \l_tmpa_tl
    \group_end:
  \fi
}
\ExplSyntaxOff
\providecommand{\tr}{\operatorname{tr}}
\newcommand{\tnleanIssueBase}{https://github.com/LionSR/TNLean/issues/}
\newcommand{\tnleanPrBase}{https://github.com/LionSR/TNLean/pull/}
\newcommand{\tnleanDocBase}{https://sirui-lu.com/TNLean/docs/}
\newcommand{\ghissue}[1]{\href{\tnleanIssueBase#1}{GitHub issue \##1}}
\newcommand{\ghpr}[1]{\href{\tnleanPrBase#1}{PR \##1}}
\newcommand{\doclink}[1]{\href{\tnleanDocBase#1.html}{\path{#1}}}

% Dirac notation and the identity operator, as in blueprint/src/macros/common.tex.
% \providecommand keeps note-local definitions authoritative.
\usepackage{dsfont}
\providecommand{\ket}[1]{|#1\rangle}
\providecommand{\bra}[1]{\langle #1|}
\providecommand{\braket}[2]{\langle #1 | #2 \rangle}
\providecommand{\ketbra}[2]{|#1\rangle\!\langle #2|}
\providecommand{\Id}{\mathds{1}}
\providecommand{\one}{\mathds{1}}
\providecommand{\C}{\mathbb{C}}
\providecommand{\MN}[1]{M_{#1}(\C)}
\providecommand{\MD}{M_D(\C)}

% Verdict marker: \gapnote{<kind>}{<status>}. Kind is the policy
% classification (clarification, local-correction, scope-restriction,
% unfaithful, false-source, open-gap); status is open, wip (elimination
% underway), resolved, or historical. Machine-read by the site index;
% prints a verdict line.
\newcommand{\gapnote}[2]{\par\noindent\textbf{Verdict:} #1 (#2).\par}


% Fallbacks keep this note compilable when a different command.tex is found first.
\providecommand{\leanid}[1]{\textsf{\detokenize{#1}}}
\providecommand{\gapnote}[2]{\par\noindent\textbf{Verdict:} #1 (#2).\par}

\title{Nontrivial Chains in the Viable Leading Step of Landau, Vazirani and Vidick}
\author{}
\date{2026-10-27}

\begin{document}
\maketitle
\gapnote{scope-restriction}{open}

\section{The source statement}

The size-trimming step of Landau, Vazirani and Vidick~\cite{Landau2015Polynomial}
(arXiv:1307.5143v1, \path{algo53.tex}, equation \path{eq:conv}, lines~472--476,
and Lemma \path{l:leading}, lines~493--505) is stated for every chain of $n$
qudits of local dimension $d$.

\section{What is formalized}

Write $L$ for the number of sites and $q$ for the local dimension. The
viable-witness and leading-vector theorems assume $L\ge1$ and $q\ge2$. These
restrictions are separate from the constant bond dimension, which is
recorded in the note on the constant bond dimension. They give the estimate
$\Delta\le4$ for a range-one Hamiltonian with interaction norm at most one,
which in turn controls $c=(\Delta/169)^2$ and the energy error in the
leading-vector argument.

\section{Elimination}

For $q=1$, the chain Hilbert space is one-dimensional and the state and
energy conclusions can be proved directly, without a gap upper bound; such a
statement must also specify a bond cutoff of at least one. The empty chain
requires a separate formulation of the quantities $c/L$ and $c/(2L)$. These
separate statements remain to be proved as part of the viable-set
construction tracked in
issue~\href{https://github.com/LionSR/TNLean/issues/9820}{9820}.

\bibliographystyle{alpha}
\bibliography{references}

\end{document}
